\subsection{2.2.2 Die Gram-Matrix als emergente Geometrie}

Nachdem im vorherigen Abschnitt der relationale Kernel

\[
R_{ij}
\]

als fundamentale Struktur des OPP eingeführt wurde, stellt sich unmittelbar die nächste Frage:

\begin{center}
\emph{Wie entsteht aus einer reinen Relation überhaupt Geometrie?}
\end{center}

Genau hier beginnt der eigentliche mathematische Kern der PST.

%%%%%%%%%%%%%%%%%%%%%%%%%%%%%%%%%%%%%%%%%%%%%%%%%%%%%%%%%%%%%%%

\subsubsection*{Vom Beziehungsnetz zur Geometrie}

Die Matrix

\[
R=(R_{ij})
\]

enthält ausschließlich paarweise Relationen.

Sie besitzt zunächst keinerlei:

\begin{itemize}
\item Koordinaten,
\item Abstände,
\item Winkel,
\item Raumdimension,
\item Metrik.
\end{itemize}

Sie beschreibt lediglich,

\[
\text{„Wie stark stehen zwei Attraktoren miteinander in Beziehung?“}
\]

Die PST behauptet deshalb:

\begin{center}
Geometrie existiert im OPP nicht fundamental.

Sie entsteht erst aus den Relationen.
\end{center}

%%%%%%%%%%%%%%%%%%%%%%%%%%%%%%%%%%%%%%%%%%%%%%%%%%%%%%%%%%%%%%%

\subsubsection*{Warum eine Gram-Matrix?}

In der klassischen Geometrie besitzt man Punkte

\[
x_i\in\mathbb R^n
\]

und definiert daraus

\[
G_{ij}=x_i\cdot x_j.
\]

Die Gram-Matrix enthält sämtliche Winkel- und Abstandsinformationen.

Die PST besitzt jedoch keine Punkte.

Deshalb muss eine Gram-Matrix ausschließlich aus Relationen konstruiert werden.

%%%%%%%%%%%%%%%%%%%%%%%%%%%%%%%%%%%%%%%%%%%%%%%%%%%%%%%%%%%%%%%

\subsubsection*{Die PST-Konstruktion}

In sämtlichen numerischen Experimenten wurde dieselbe Konstruktion verwendet.

Sei

\[
R_{ij}
\]

die relationale Stärke.

Dann wird bezüglich eines Referenzattraktors

\[
0
\]

definiert

\[
\boxed{
K_{ij}
=
R_{i0}
+
R_{j0}
-
R_{ij}.
}
\]

Diese Matrix nennen wir

\[
\boxed{\mathcal K}
\]

oder einfach

\[
K.
\]

Sie bildet den zentralen mathematischen Operator der bisherigen PST-Simulationen.

%%%%%%%%%%%%%%%%%%%%%%%%%%%%%%%%%%%%%%%%%%%%%%%%%%%%%%%%%%%%%%%

\subsubsection*{Warum genau diese Formel?}

Die Konstruktion erinnert unmittelbar an die klassische Metrikrekonstruktion.

Für euklidische Punkte gilt

\[
||x_i-x_j||^2
=
||x_i||^2
+
||x_j||^2
-
2x_i\cdot x_j.
\]

Die PST ersetzt

\[
||x_i||^2
\]

durch

\[
R_{i0},
\]

also durch die Relation zum Referenzattraktor.

Damit wird

\[
K
\]

zum Analogon einer Gram-Matrix.

Wichtig:

Es werden keinerlei Koordinaten vorausgesetzt.

Alles entsteht ausschließlich aus den Relationen.

%%%%%%%%%%%%%%%%%%%%%%%%%%%%%%%%%%%%%%%%%%%%%%%%%%%%%%%%%%%%%%%

\subsubsection*{Interpretation}

Man kann die Matrix

\[
K
\]

auf verschiedene Arten interpretieren.

\paragraph{1. Informationsgeometrie}

Jeder Eintrag beschreibt,

wie ähnlich zwei Attraktoren bezüglich des gesamten relationalen Netzwerkes sind.

\paragraph{2. Emergentes Skalarprodukt}

Formal verhält sich

\[
K
\]

wie ein Skalarprodukt.

Deshalb lassen sich Eigenwerte und Eigenvektoren interpretieren wie Raumachsen.

\paragraph{3. Projektionsoperator}

Die später entstehenden Raumdimensionen sind nichts anderes als Eigenrichtungen dieser Matrix.

%%%%%%%%%%%%%%%%%%%%%%%%%%%%%%%%%%%%%%%%%%%%%%%%%%%%%%%%%%%%%%%

\subsubsection*{Eigenschaften}

Die Matrix besitzt mehrere wichtige Eigenschaften.

\paragraph{Symmetrie}

Da

\[
R_{ij}=R_{ji}
\]

gilt,

ist unmittelbar

\[
K_{ij}=K_{ji}.
\]

Somit

\[
K
=
K^T.
\]

%%%%%%%%%%%%%%%%%%%%%%%%%%%%%%%%%%%%%%%%%%%%%%%%%%%%%%%%%%%%%%%

\paragraph{Reelles Spektrum}

Da

\[
K
\]

symmetrisch ist,

besitzt sie ausschließlich reelle Eigenwerte

\[
\lambda_i.
\]

Dies ist entscheidend für sämtliche Spektraluntersuchungen.

%%%%%%%%%%%%%%%%%%%%%%%%%%%%%%%%%%%%%%%%%%%%%%%%%%%%%%%%%%%%%%%

\paragraph{Eigenzerlegung}

Es gilt

\[
K
=
V
\Lambda
V^T.
\]

Dabei ist

\[
\Lambda
=
\operatorname{diag}
(\lambda_1,\ldots,\lambda_n).
\]

Die Eigenvektoren bilden die möglichen Projektionsrichtungen.

%%%%%%%%%%%%%%%%%%%%%%%%%%%%%%%%%%%%%%%%%%%%%%%%%%%%%%%%%%%%%%%

\subsubsection*{Die effektive Dimension}

Aus den Eigenwerten wurde in nahezu allen Tests dieselbe Größe berechnet.

Die sogenannte effektive Dimension

\[
d_{\mathrm{eff}}
\]

lautet

\[
\boxed{
d_{\mathrm{eff}}
=
\frac{\left(\sum_i |\lambda_i|\right)^2}
{\sum_i\lambda_i^2}.
}
\]

Diese Größe besitzt mehrere bemerkenswerte Eigenschaften.

%%%%%%%%%%%%%%%%%%%%%%%%%%%%%%%%%%%%%%%%%%%%%%%%%%%%%%%%%%%%%%%

\paragraph{Beispiele}

Besitzen alle Eigenwerte denselben Betrag,

\[
\lambda_i=1
\]

für

\[
i=1,\ldots,n,
\]

so ergibt sich

\[
d_{\mathrm{eff}}
=
n.
\]

%%%%%%%%%%%%%%%%%%%%%%%%%%%%%%%%%%%%%%%%%%%%%%%%%%%%%%%%%%%%%%%

Ist dagegen nur ein Eigenwert ungleich Null,

\[
\lambda_1\neq0,
\qquad
\lambda_{i>1}=0,
\]

dann folgt

\[
d_{\mathrm{eff}}
=
1.
\]

%%%%%%%%%%%%%%%%%%%%%%%%%%%%%%%%%%%%%%%%%%%%%%%%%%%%%%%%%%%%%%%

Die Größe misst also nicht die Anzahl mathematischer Dimensionen,

sondern

\begin{center}
die Anzahl effektiv aktiver Freiheitsgrade.
\end{center}

%%%%%%%%%%%%%%%%%%%%%%%%%%%%%%%%%%%%%%%%%%%%%%%%%%%%%%%%%%%%%%%

\subsubsection*{Warum genau diese Definition?}

Diese Definition entspricht der sogenannten Participation Ratio,

die u.a. verwendet wird in

\begin{itemize}
\item Random-Matrix-Theorie,
\item Quantenchaos,
\item Anderson-Lokalisierung,
\item Netzwerktheorie,
\item Informationstheorie.
\end{itemize}

Sie misst,

wie viele Eigenmoden tatsächlich zum Gesamtsystem beitragen.

%%%%%%%%%%%%%%%%%%%%%%%%%%%%%%%%%%%%%%%%%%%%%%%%%%%%%%%%%%%%%%%

\subsubsection*{Physikalische Interpretation}

In der PST besitzt

\[
d_{\mathrm{eff}}
\]

eine wesentlich tiefere Bedeutung.

Sie beschreibt nicht

\[
\text{„Wie viele Dimensionen existieren?“}
\]

sondern

\[
\boxed{
\text{„Wie viele Freiheitsgrade werden durch die Projektion tatsächlich sichtbar?“}
}
\]

Damit wird

\[
d_{\mathrm{eff}}
\]

zur mathematischen Brücke zwischen

\[
\text{OPP}
\]

und

\[
\text{emergenter Raumzeit}.
\]

%%%%%%%%%%%%%%%%%%%%%%%%%%%%%%%%%%%%%%%%%%%%%%%%%%%%%%%%%%%%%%%

\subsubsection*{Konzeptionelle Bedeutung}

Diese Sichtweise verändert den gesamten geometrischen Ausgangspunkt.

Die klassische Physik beginnt mit

\[
\text{Raum}
\rightarrow
\text{Punkte}
\rightarrow
\text{Abstände}.
\]

Die PST beginnt stattdessen mit

\[
\boxed{
\text{Relation}
\rightarrow
\text{Gram-Matrix}
\rightarrow
\text{Eigenstruktur}
\rightarrow
\text{Geometrie}.
}
\]

Raum ist somit kein Fundament,

sondern ein Spektrum.

%%%%%%%%%%%%%%%%%%%%%%%%%%%%%%%%%%%%%%%%%%%%%%%%%%%%%%%%%%%%%%%

\subsubsection*{Folgerung}

Die Gram-Matrix stellt den ersten Schritt dar,

bei dem reine Information in geometrische Struktur übergeht.

Sie besitzt keinerlei klassische Rauminterpretation,

enthält jedoch bereits sämtliche Informationen,

aus denen eine Raumzeitprojektion entstehen kann.

Alle späteren numerischen Untersuchungen — Optimierungen, Spektralkollaps, Diffusionsprojektionen, Kernel-PCA und Stabilitätsfunktionale — operieren ausschließlich auf dieser Matrix.

Damit bildet

\[
\boxed{\mathcal K}
\]

den eigentlichen mathematischen Mittelpunkt der bisherigen PST.

